\chapter{Путеводитель по параллелизму для путешествующих автостопом}
\label{the-hitchhikers-guide-to-concurrency}
Далеко\--далеко в закоулках нефешенебельного начала 21\--го века, которого даже нет на карте, находится маленькое подмножество человеческих знаний.

В этом подмножестве заключена совершенно невзрачная маленькая дисциплина, чья Фон\--Неймановская архитектура столь примитивна, что согласно ей ОПЗ калькуляторы считаются чем\--то выдающимся.

У этой дисциплины есть, а точнее была \--- проблема: большинство людей, изучающих её, были почти всегда недовольны, создавая параллельное ПО.
Предлагалось множество решений этой проблемы, но чаще всего они были связаны с манипуляцией маленькими логическими блоками, которые назывались локами, мутексами и всякими другими именами, что несколько странно, поскольку этим самым блокам параллелизм был совершенно не нужен.

Так проблема и оставалась нерешённой: одни люди были жадными, другие жалкими, и даже ОПЗ калькуляторы им не помогали.

Кое\--кто был убеждён, что людям не стоило добавлять параллелизм в языки программирования.
И что программы вообще не должны были выходить из их первоначального потока.
\\
\colorbox{lgray}
{
\begin{minipage}{1.0\linewidth}
    \textbf{Замечание:} неплохое развлечение \--- писать пародии на <<Путеводитель по галактике для путешествующих автостопом>>.
    Если вам ещё не попадалась эта книга, то обязательно её прочитайте.
    Она стоит того!
\end{minipage}
}
\section{Без паники}
\label{dont-panic}
\begin{wrapfigure}{l}{0.35\linewidth}
    \includegraphics[width=1\linewidth]{fat-guy.png}
\end{wrapfigure}
Привет.
Сегодня (или в любой из дней, когда вы читаете эти строки, даже завтра) я хочу рассказать вам о параллельном Erlang.
Скорее всего вы уже читали о параллелизме, или когда\--либо сталкивались с ним.
Может быть, вы интересуетесь истоками программирования для множества ядер, или читаете эту книгу, наслушавшись болтовни, которой в наше время окружён параллелизм.

Впрочем, хочу вас предупредить, что в этой главе основной упор сделан на теорию.
Если у вас болит голова, или вы питаете отвращение к истории языков программирования, или вам просто хотелось попрограммировать, то вам лучше перейти к ~\ref{thanks-for-all-the-fish}~концу главы, или перейти к следующей (там как раз освещается более практическая сторона вопроса).

Я уже объяснил во введении к книге, что параллелизм в Erlang основан на передаче сообщений и модели акторов.
Мой пример рассказывал о людях, общение которых происходит исключительно при помощи писем.
Чуть позже я ещё вернусь к этой модели, а сейчас нам необходимо первым делом обозначить разницу между \emph{конкурентностью} и \emph{параллелизмом}.

Оба слова во многих ситуациях имеют одно и то же значение.
Но в контексте Erlang они часто относятся к двум разным концепциям.
По мнению большинства эрлангистов понятие конкурентности описывает несколько акторов, которые исполняются независимо друг от друга, но их исполнение не обязательно происходит в один тот же момент.
Параллелизм же означает, что несколько акторов исполняются одновременно.
Взгляды различных областей computer science на правильность этих определений могут не совпадать, но в этом руководстве я буду пользоваться именно такими определениями.
Не удивляйтесь, если в других источниках вы увидите, как кто\--то употребляет те же самые термины для обозначения других понятий.

Конкурентность была в Erlang с самого начала, даже в восьмедисятые, когда всё запускалось на одноядерном процессоре.
Каждому процессу Erlang для исполнения отводился свой собственный временной отрезок, совсем как в эру десктопных приложений, которая предшествовала появлению многоядерных систем.

Уже в то время можно было, в принципе, реализовать параллелизм.
Для этого нам бы потребовался второй компьютер, который бы исполнял код и обменивался информацией с первым.
Но даже такая система могла бы исполнять параллельно всего лишь два актора.
Современные многоядерные системы позволяют реализовать параллелизм в рамках одного компьютера (некоторые промышленные чипы могут содержать десятки ядер), и Erlang использует эту возможность в полной мере.\\
\colorbox{lorange}
{
\begin{minipage}{1.0\linewidth}
    \textbf{Не забывайтесь:}\\
    Важно понимать разницу между конкурентностью и параллелизмом.
    Многие программисты верят, что Erlang был готов к использованию на многоядерных компьютерах задолго до того, как это произошло в действительности.
    Erlang стал использовать истинную \href{http://en.wikipedia.org/wiki/Symmetric\_multiprocessing}{симметричную мультипроцессорность} (symmetric multiprocessing) лишь в середине двухтысячных, а большая часть реализации была завершена в релизе R13B в 2009 году.
    До этого часто приходилось отключать SMP, чтобы избежать потерь производительности.
    На многоядерном компьютере без SMP можно получить параллелизм, если запустить одновременно несколько экземпляров виртуальной машины.\\
    \\
    Интересно отметить, что для внесения истинного параллелизма в язык, не потребовалось производить какие\--либо концептуальные изменения на языковом уровне.
    Благодаря тому, что конкурентность в Erlang строится вокруг изолированных процессов, все изменения были сделаны внутри ВМ, подальше от глаз обычного программиста.
\end{minipage}
}
